% ============================================================
%  Lernzettel Template (my_dhbw Stil, analysis.tex)
%  Eine Spalte, mit TOC, accent-Theme.
%  Renderer: pdflatex (zweimal)
% ============================================================
\documentclass[11pt, a4paper]{article}

\usepackage[ngerman]{babel}
\usepackage{fontspec}
\setmainfont[
  Path=/usr/share/fonts/TTF/,
  Extension=.ttf,
  BoldFont=DejaVuSans-Bold,
  ItalicFont=DejaVuSans-Oblique,
  BoldItalicFont=DejaVuSans-BoldOblique
]{DejaVuSans}
\setsansfont[
  Path=/usr/share/fonts/TTF/,
  Extension=.ttf,
  BoldFont=DejaVuSans-Bold,
  ItalicFont=DejaVuSans-Oblique,
  BoldItalicFont=DejaVuSans-BoldOblique
]{DejaVuSans}
\setmonofont[Path=/usr/share/fonts/TTF/,Extension=.ttf]{DejaVuSansMono}
\usepackage[top=2cm, bottom=2cm, left=2cm, right=2cm, footskip=1cm]{geometry}
\usepackage{amsmath, amssymb}
\usepackage{xcolor}
\usepackage{titlesec}
\usepackage{titletoc}
\usepackage{enumitem}
\usepackage{fancyhdr}
\usepackage{lastpage}
\usepackage{microtype}
\usepackage{parskip}

\usepackage{booktabs}
\usepackage{array}
\usepackage{longtable}
\usepackage{tabularx}
\usepackage{ltablex}
\keepXColumns
\usepackage{colortbl}
\usepackage[most,breakable]{tcolorbox}
\usepackage{listings}
\usepackage[normalem]{ulem}
\usepackage{pdflscape}
\usepackage[colorlinks=true, linkcolor=accent, urlcolor=accent]{hyperref}

\renewcommand{\familydefault}{\sfdefault}

% ── Theme (my_dhbw accent) ────────────────────────────────────
\definecolor{accent}{RGB}{0, 60, 113}
\definecolor{primaryblue}{RGB}{0, 60, 113}
\definecolor{codebg}{RGB}{245, 245, 245}
\definecolor{figurebg}{RGB}{232, 241, 255}
\definecolor{tablehintbg}{RGB}{240, 255, 240}
\definecolor{solutionbg}{RGB}{242, 255, 242}
\definecolor{rulegray}{RGB}{180, 180, 180}
\definecolor{tableheadbg}{RGB}{0, 60, 113}

% ── Header/Footer ─────────────────────────────────────────────
\pagestyle{fancy}
\fancyhf{}
\renewcommand{\headrulewidth}{0pt}
\renewcommand{\footrulewidth}{0.4pt}
\renewcommand{\sectionmark}[1]{\markboth{#1}{}}
\fancyfoot[L]{\footnotesize\color{accent}\nouppercase{\leftmark}}
\fancyfoot[R]{\footnotesize\color{accent}\thepage\,/\,\pageref{LastPage}}
\fancypagestyle{plain}{%
  \fancyhf{}%
  \renewcommand{\headrulewidth}{0pt}%
  \renewcommand{\footrulewidth}{0.4pt}%
  \fancyfoot[L]{\footnotesize\color{accent}\nouppercase{\leftmark}}%
  \fancyfoot[R]{\footnotesize\color{accent}\thepage\,/\,\pageref{LastPage}}%
}

% ── Typografie ────────────────────────────────────────────────
\titleformat{\section}
    {\color{accent}\large\bfseries}{}{0pt}{}
    [\vspace{-4pt}\color{accent}\rule{\linewidth}{2pt}]
\titleformat{\subsection}
    {\normalsize\bfseries\color{accent!85!black}}{}{0pt}{}
    [\vspace{-3pt}\color{accent!40}\rule{\linewidth}{0.4pt}]
\titleformat{\subsubsection}
    {\small\bfseries\color{accent!70!black}}{}{0pt}{}{}
\titlespacing*{\section}{0pt}{12pt}{4pt}
\titlespacing*{\subsection}{0pt}{8pt}{2pt}
\titlespacing*{\subsubsection}{0pt}{6pt}{1pt}

\setlength{\parindent}{0pt}
\setlength{\parskip}{4pt}
\renewcommand{\arraystretch}{1.25}
\setlist[itemize]{noitemsep, topsep=3pt, parsep=0pt, leftmargin=1.4em}
\setlist[enumerate]{noitemsep, topsep=3pt, parsep=0pt, leftmargin=1.6em}
\setlist[itemize,2]{label=\textendash, leftmargin=1.4em}

% ── Boxen ──────────────────────────────────────────────────────
\newtcolorbox{figurebox}[1]{
  colback=figurebg, colframe=accent!50,
  title={Abbildung: #1}, fonttitle=\small\bfseries,
  breakable, before skip=6pt, after skip=6pt
}
\newtcolorbox{tablehintbox}[1]{
  colback=tablehintbg, colframe=green!50!black!50,
  title={Tabelle: #1}, fonttitle=\small\bfseries,
  breakable, before skip=6pt, after skip=6pt
}
\newtcolorbox{solutionbox}[1]{
  colback=solutionbg, colframe=green!60!black!60,
  title={#1}, fonttitle=\small\bfseries,
  breakable, before skip=6pt, after skip=6pt
}

\lstset{
  basicstyle=\ttfamily\small,
  backgroundcolor=\color{codebg},
  breaklines=true,
  frame=single, framerule=0pt,
  xleftmargin=4pt, xrightmargin=4pt,
  aboveskip=6pt, belowskip=6pt,
  keepspaces=true
}

% ── TOC-Look ──────────────────────────────────────────────────
\renewcommand{\contentsname}{\color{accent}Inhalt}

\providecommand{\nichttitle}{Lernzettel}
\providecommand{\nichtsubtitle}{}

\renewcommand{\nichttitle}{Digitale Betriebswirtschaftslehre}
\renewcommand{\nichtsubtitle}{Lernzettel}


\begin{document}

\begin{center}
{\bfseries\Large\color{accent}\nichttitle}\\[2pt]
\color{accent}\rule{0.35\linewidth}{1pt}\color{black}
\end{center}
\vspace{2pt}

{\small\tableofcontents}
\newpage

% ── BODY ──────────────────────────────────────────────────────

\section{Digitale Betriebswirtschaftslehre — Gesamtzettel}


\subsection{1. Grundlagen der BWL}

\textbf{Wirtschaften}: Knappe Güter geplant einsetzen, sodass Bedürfnisbefriedigung möglichst vorteilhaft erfolgt.

\begin{itemize}
  \item Bedürfnisse + begrenzte Ressourcen → Knappheit → wirtschaftliches Handeln
\end{itemize}

\textbf{Wertschöpfungssystem}: Input (Produktionsfaktoren) → Transformation → Output (Produkte, Dienstleistungen, Wert).

\begin{itemize}
  \item Wertschöpfung = Output − Input
  \item Modell: Porters Wertschöpfungskette
\end{itemize}


\subsection{2. Ansätze der BWL}


{\small
\begin{tabularx}{\linewidth}{XXXXXX}
\hline
\rowcolor{tableheadbg}
{\color{white}\textbf{Ansatz}} & {\color{white}\textbf{Vertreter}} & {\color{white}\textbf{Blick auf Unternehmen}} & {\color{white}\textbf{Zentrale Elemente}} & {\color{white}\textbf{Kernfrage}} & {\color{white}\textbf{Zielgröße}} \\[2pt]
\hline
\endhead
\hline
\endfoot
Produktionsorientiert & Erich Gutenberg & Produktion & Arbeit, Kapital, Material & Wie effizient produzieren? & Produktivität (Output/Input) \\
Entscheidungsorientiert & Edmund Heinen & Entscheidungen & Daten, Alternativen, Modelle & Wie gute Entscheidungen treffen? & Qualität von Entscheidungen \\
Systemorientiert & Hans Ulrich & Gesamtsystem & Elemente, Umwelt, Rückkopplungen & Wie hängt alles zusammen? & Stabilität, Überlebensfähigkeit \\
Institutionenökonomisch & Ronald Coase & Regeln \& Verträge & Akteure, Verträge, Anreize & Wie effizient koordinieren? & Minimierung Transaktionskosten \\
\end{tabularx}
}


\textbf{Systemorientierter Ansatz}: Unternehmen = komplexe, offene Systeme.

\begin{itemize}
  \item Vernetzte Elemente (Abteilungen, Prozesse, Menschen), Interaktion mit Umwelt (Markt, Technologie, Gesellschaft), Rückkopplungen/Regelkreise
\end{itemize}

\textbf{Institutionenökonomischer Ansatz}: Analyse Transaktionskosten, Informationsasymmetrien, Principal-Agent-Probleme.

\begin{itemize}
  \item Idee: Unternehmen organisieren Transaktionen effizienter als der Markt
\end{itemize}


\subsection{3. Produktionsfaktoren}

\textbf{Klassisch (Gutenberg)}:

\begin{itemize}
  \item \textbf{Arbeit}: menschliche Arbeitsleistung
  \item \textbf{Betriebsmittel (Kapital)}: Maschinen, Anlagen, IT-Systeme
  \item \textbf{Werkstoffe}: Rohstoffe, Materialien
\end{itemize}

\textbf{Daten als moderner Produktionsfaktor}: tragen zur Wertschöpfung bei via Sammlung → Verarbeitung → Entscheidung.

\begin{itemize}
  \item Ermöglichen bessere Prognosen, effizientere Prozesse, neue Geschäftsmodelle
  \item Wandel: Input → Produktion → Output  ⇒  Daten → Analyse → Entscheidung → Wertschöpfung
\end{itemize}


\subsection{4. Principal-Agent-Theorie}

\textbf{Principal-Agent-Theorie}: Trennung von Eigentum und Kontrolle führt zu Interessenkonflikten.

\begin{itemize}
  \item \textbf{Principal}: Eigentümer (delegiert)
  \item \textbf{Agent}: Manager (entscheidet)
  \item \textbf{Informationsasymmetrie}: Principal hat weniger Infos über Handeln des Agents
  \item Risiko: Agent verfolgt eigene Interessen (Macht, kurzfristige Boni)
  \item Lösung: Anreiz-/Kontrollsysteme — leistungsabhängige Vergütung, KPIs, Monitoring
\end{itemize}


\subsection{5. Shareholder vs. Stakeholder}


{\small
\begin{tabularx}{\linewidth}{XXX}
\hline
\rowcolor{tableheadbg}
{\color{white}\textbf{Merkmal}} & {\color{white}\textbf{Stakeholder-Ansatz}} & {\color{white}\textbf{Shareholder-Ansatz}} \\[2pt]
\hline
\endhead
\hline
\endfoot
Fokus & alle Anspruchsgruppen & Eigentümer (Aktionäre) \\
Oberstes Ziel & Interessen aller Gruppen & Maximierung Unternehmenswert \\
Erfolgsmaßstab & Maximierung Differenz Nutzen − Kosten aller & Maximierung zukünftiger Zahlungen an Eigentümer \\
Beurteilung & nicht operational, pluralistisch & operational, monistisch \\
\end{tabularx}
}



\subsubsection{Stakeholder-Übersicht}


{\small
\begin{tabularx}{\linewidth}{XXX}
\hline
\rowcolor{tableheadbg}
{\color{white}\textbf{Stakeholder}} & {\color{white}\textbf{Typ}} & {\color{white}\textbf{Anspruch}} \\[2pt]
\hline
\endhead
\hline
\endfoot
Eigenkapitalgeber & intern & Kapitalmehrung (Ausschüttung, Kapitalzuwachs) \\
Arbeitnehmer & intern & Entlohnung, Arbeitsbedingungen, Sicherheit \\
Management & intern & Gehalt, Einfluss, Status \\
Fremdkapitalgeber & extern & Tilgung und Verzinsung \\
Kunden & extern & Wettbewerbsfähige Güter, Qualität \\
Lieferanten & extern & Zuverlässige Zahlungen, Langfristigkeit \\
Staat/Gesellschaft & extern & Steuern, Recht, Umweltschutz \\
\end{tabularx}
}



\subsection{6. Unternehmensziele}

\textbf{Unternehmensziele}: Maßstäbe, an denen zukünftiges Handeln ausgerichtet und bewertet wird.

\begin{itemize}
  \item Funktionen: Orientierung, Entscheidungskriterium, Handlungssteuerung, Koordination, Legitimation/Konfliktlösung, Bewertung/Kontrolle
\end{itemize}


\subsubsection{Zielinhalt}


{\small
\begin{tabularx}{\linewidth}{XXX}
\hline
\rowcolor{tableheadbg}
{\color{white}\textbf{}} & {\color{white}\textbf{Sachziele}} & {\color{white}\textbf{Formalziele (Erfolgsziele)}} \\[2pt]
\hline
\endhead
\hline
\endfoot
Bezug & Konkretes Handeln in betrieblichen Funktionen & Übergeordnete Ziele, an denen Sachziele auszurichten sind \\
Beispiele & Produkte, Qualität, Marktposition & Gewinn, Rentabilität, Liquidität \\
\end{tabularx}
}


\textbf{Sachziel-Kategorien}: Leistungsziele, Finanzziele, Führungs-/Organisationsziele, soziale/ökologische Ziele.



\subsubsection{Zielausmaß}
\begin{itemize}
  \item Wie stark soll Ziel erreicht werden, wie messen? (z. B. Gewinn ≥ 10 \% vs. Gewinn maximieren)
  \item Messbarkeit: kardinal (€), ordinal (gut/schlecht), nominal (erreicht/nicht)
\end{itemize}


\subsubsection{Zeitbezug}
\begin{itemize}
  \item kurzfristig ≤ 1 Jahr | mittelfristig 1–5 Jahre | langfristig > 5 Jahre
\end{itemize}


\subsubsection{Rangordnung}
\begin{itemize}
  \item Hauptziele vs. Nebenziele; Ober-, Zwischen-, Unterziele → Zielhierarchie
\end{itemize}


\subsubsection{Zielbeziehungen}


{\small
\begin{tabularx}{\linewidth}{XX}
\hline
\rowcolor{tableheadbg}
{\color{white}\textbf{Typ}} & {\color{white}\textbf{Bedeutung}} \\[2pt]
\hline
\endhead
\hline
\endfoot
Komplementär & Erreichen von Ziel A fördert Ziel B \\
Konkurrierend (konfliktär) & Erreichen von Ziel A behindert Ziel B \\
Indifferent (neutral) & Ziele beeinflussen sich nicht \\
\end{tabularx}
}


\textbf{Beispielpaare}:

\begin{itemize}
  \item Personalisierter Service ↔ Reduktion stationäres Service-Netz → konkurrierend
  \item KI-Automatisierung Fertigung ↔ Reduktion Ausschuss → komplementär
  \item Flexibilisierung Arbeitszeiten ↔ Steigerung Vertrieb → indifferent
  \item Just-in-Time ↔ Reduktion Lagerkosten → komplementär
  \item Produktivität (digitale Tools) ↔ Verschiebung Investitionen → konkurrierend
  \item Stilllegung Produktionskapazitäten ↔ Absatzsteigerung bei konstantem Preis → konkurrierend
\end{itemize}


\subsubsection{SMART}


{\small
\begin{tabularx}{\linewidth}{XXX}
\hline
\rowcolor{tableheadbg}
{\color{white}\textbf{Buchstabe}} & {\color{white}\textbf{Bedeutung}} & {\color{white}\textbf{Erklärung}} \\[2pt]
\hline
\endhead
\hline
\endfoot
S & Spezifisch & konkret, unmissverständlich \\
M & Messbar & quantitativ/qualitativ beurteilbar \\
A & Attraktiv & angemessen, motivierend \\
R & Realistisch & mit Ressourcen erreichbar \\
T & Terminiert & bis zu bestimmtem Zeitpunkt \\
\end{tabularx}
}



\subsection{7. Formalziele — Kennzahlen}

\textbf{Produktivität}: mengenmäßiges Verhältnis Output/Input — misst Ergiebigkeit.

\begin{itemize}
  \item Produktivität = Output-Menge / Input-Menge
\end{itemize}

\textbf{Wirtschaftlichkeit}: wertmäßiges Verhältnis — misst Effizienz in Geldeinheiten.

\begin{itemize}
  \item Wirtschaftlichkeit = Ertrag (€) / Aufwand (€)
\end{itemize}

\textbf{Gewinn} = Ertrag − Aufwand


\textbf{Rentabilität} = (Gewinn / durchschnittlich eingesetztes Kapital) × 100 — relativer Erfolg zum Kapital



\subsubsection{Beispiel: Schraubenproduktion}
\begin{itemize}
  \item Output: 1.000 Schrauben | Input: 10 kg Draht | Drahtpreis: 2 €/kg | Verkaufspreis: 0,02 €/Schraube
  \item Produktivität = 1.000 / 10 = 100 Schrauben/kg
  \item Aufwand = 10 · 2 = 20 € | Ertrag = 1.000 · 0,02 = 20 €
  \item Wirtschaftlichkeit = 20 / 20 = 1,0
  \item Bei +10 \%: Produktivität = 110 Schrauben/kg; Wirtschaftlichkeit = 1,1
\end{itemize}


\subsubsection{Hebel zur Verbesserung}
\begin{itemize}
  \item Mengen: Output erhöhen oder Input senken
  \item Preise: Verkaufspreis erhöhen oder Einkaufspreis senken
\end{itemize}


\subsection{8. Ökonomisches Prinzip}


{\small
\begin{tabularx}{\linewidth}{XX}
\hline
\rowcolor{tableheadbg}
{\color{white}\textbf{Prinzip}} & {\color{white}\textbf{Logik}} \\[2pt]
\hline
\endhead
\hline
\endfoot
Maximum-Prinzip & Gegebener Input → maximaler Output \\
Minimum-Prinzip & Gegebener Output → minimaler Input \\
Optimum-Prinzip & möglichst günstiges Verhältnis Output/Input \\
\end{tabularx}
}



\subsection{9. Effizienz vs. Effektivität}


{\small
\begin{tabularx}{\linewidth}{XXX}
\hline
\rowcolor{tableheadbg}
{\color{white}\textbf{}} & {\color{white}\textbf{Effizienz}} & {\color{white}\textbf{Effektivität}} \\[2pt]
\hline
\endhead
\hline
\endfoot
Bedeutung & Verhältnis Leistung/Ressourceneinsatz & Grad der Zielerreichung \\
Frage & „Werden die Dinge richtig gemacht?" & „Werden die richtigen Dinge gemacht?" \\
Begriff & Leistungsfähigkeit & Leistungswirksamkeit \\
\end{tabularx}
}



\subsection{10. Mindestbedingungen wirtschaftlichen Handelns}

\begin{itemize}
  \item \textbf{Langfristig (Erfolg)}: Erträge ≥ Aufwendungen → keine Dauerverluste
  \item \textbf{Kurzfristig (Liquidität)}: Einzahlungen ≥ Auszahlungen → Zahlungsfähigkeit
\end{itemize}


\subsection{11. Markt — Grundlagen}

\textbf{Markt}: Ort, an dem Angebot und Nachfrage aufeinandertreffen — koordiniert wirtschaftliche Aktivitäten durch Austausch Güter/Dienstleistungen ↔ Geld.

\begin{itemize}
  \item Marktteilnehmer: Anbieter (Unternehmen), Nachfrager (Kunden)
\end{itemize}


\subsubsection{Grundmodell}
\begin{itemize}
  \item \textbf{Nachfrage}: Zahlungsbereitschaft der Kunden; Preis ↑ → Nachfrage ↓
  \item \textbf{Angebot}: Verkaufsbereitschaft der Unternehmen; Preis ↑ → Angebot ↑
\end{itemize}


\subsubsection{Grenzen des Modells}
\begin{itemize}
  \item Informationsasymmetrien, Markteintrittsbarrieren, Marktmacht einzelner Anbieter
\end{itemize}


\subsection{12. Marktgleichgewicht}

\textbf{Gleichgewicht}: Preis bei dem QD = QS — markträumend, ohne Über-/Unterversorgung.

\begin{itemize}
  \item \textbf{Gleichgewichtspreis}: Preis ohne Über-/Unterversorgung
  \item \textbf{Gleichgewichtsmenge}: gehandelte Menge im Markt
  \item QD > QS → Nachfrageüberschuss (Preis zu niedrig) → Preis steigt
  \item QS > QD → Angebotsüberschuss (Preis zu hoch) → Preis sinkt
\end{itemize}


\subsubsection{Berechnung — Beispiel Ladekabel}
\begin{itemize}
  \item Nachfrage: QD = 100 − 2P | Angebot: QS = 20 + 2P
  \item Gleichsetzen: 100 − 2P = 20 + 2P
  \item 80 = 4P → P* = 20 €
  \item Q* = 100 − 2·20 = 60
\end{itemize}


\subsubsection{Disequilibrium bei P = 10 €}
\begin{itemize}
  \item QD = 80 | QS = 40 | Nachfrageüberschuss = 40 → Preis steigt bis P* = 20
\end{itemize}


\subsubsection{Nachfrageverschiebung: QD = 120 − 2P}
\begin{enumerate}
  \item 120 − 2P = 20 + 2P
  \item 100 = 4P → P* = 25 €
  \item Q* = 70
  \item Rechtsverschiebung Nachfragekurve → höherer Preis, höhere Menge
  \item Managementableitung: Kapazitäten ausbauen, Preis anheben
\end{enumerate}


\subsubsection{Datenbasierte Nachfrageschätzung}
\begin{itemize}
  \item Wettbewerbsvorteil: frühe Erkennung von Nachfrageveränderungen via Verkaufszahlen, Suchanfragen → schnellere Preis-/Mengenanpassung
\end{itemize}


\subsection{13. Marktformen}


{\small
\begin{tabularx}{\linewidth}{XXXX}
\hline
\rowcolor{tableheadbg}
{\color{white}\textbf{}} & {\color{white}\textbf{viele Anbieter}} & {\color{white}\textbf{wenige Anbieter}} & {\color{white}\textbf{ein Anbieter}} \\[2pt]
\hline
\endhead
\hline
\endfoot
\textbf{viele Nachfrager} & Polypol & Oligopol & Monopol \\
\textbf{wenige Nachfrager} & Nachfrageoligopol & zweiseitiges Oligopol & beschränktes Monopol \\
\textbf{ein Nachfrager} & Nachfragemonopol & beschränktes Monopol & zweiseitiges Monopol \\
\end{tabularx}
}



\subsubsection{Vergleich}

{\small
\begin{tabularx}{\linewidth}{XXXX}
\hline
\rowcolor{tableheadbg}
{\color{white}\textbf{Merkmal}} & {\color{white}\textbf{Polypol}} & {\color{white}\textbf{Oligopol}} & {\color{white}\textbf{Monopol}} \\[2pt]
\hline
\endhead
\hline
\endfoot
Anbieter & viele & wenige & einer \\
Marktmacht & keine & mittel & hoch \\
Wettbewerb & sehr hoch & begrenzt & gering \\
\end{tabularx}
}



\subsection{14. Wettbewerb}

\begin{itemize}
  \item Unternehmen konkurrieren um Kunden → Ziel: nachhaltiger Wettbewerbsvorteil
  \item Formen: \textbf{Kostenführerschaft}, \textbf{Differenzierung}, \textbf{Fokussierung}
\end{itemize}

\textbf{Markteintrittsbarrieren}: erschweren/verhindern Eintritt neuer Anbieter.

\begin{itemize}
  \item Typische: hohe Investitionskosten, Skaleneffekte, Vertriebskanäle, Regulierung
  \item Folgen: weniger Wettbewerb, stärkere Marktmacht Bestehender
\end{itemize}


\subsection{15. Digitale Märkte \& Plattformen}

\begin{itemize}
  \item \textbf{Plattformmärkte}: vermitteln zwischen ≥ 2 Nutzergruppen (Amazon: Händler–Kunden; Uber: Fahrer–Fahrgäste; Airbnb: Vermieter–Gäste)
\end{itemize}

\textbf{Netzwerkeffekte}: mehr Nutzer → höherer Nutzen für alle → schnelles Wachstum, Marktkonzentration.


\textbf{Lock-in-Effekte}: hohe Wechselkosten/Gewohnheit binden Kunden, erschweren Anbieterwechsel.



\subsection{16. Entscheidung — Grundlagen}

\textbf{Entscheidung}: Auswahl einer Alternative aus mehreren Handlungsmöglichkeiten.



\subsubsection{Arten}


{\small
\begin{tabularx}{\linewidth}{XXX}
\hline
\rowcolor{tableheadbg}
{\color{white}\textbf{Art}} & {\color{white}\textbf{Merkmale}} & {\color{white}\textbf{Beispiele}} \\[2pt]
\hline
\endhead
\hline
\endfoot
Konstitutiv & langfristig, grundlegend & Standortwahl, Rechtsform, Geschäftsmodell \\
Laufend & operativ, wiederkehrend, kurzfristig & Preis, Marketingkampagne, Lagerbestand \\
\end{tabularx}
}



\subsubsection{Entscheidungstheorie}


{\small
\begin{tabularx}{\linewidth}{XXX}
\hline
\rowcolor{tableheadbg}
{\color{white}\textbf{Theorie}} & {\color{white}\textbf{Frage}} & {\color{white}\textbf{Merkmale}} \\[2pt]
\hline
\endhead
\hline
\endfoot
Normativ & „Wie SOLLTE entschieden werden?" & rational, modellbasiert (z. B. Erwartungswert) \\
Deskriptiv & „Wie wird TATSÄCHLICH entschieden?" & psychologisch, Biases, begrenzte Rationalität \\
\end{tabularx}
}



\subsection{17. Entscheidungsproblem}

\textbf{Entscheidungsproblem}: Strukturierung der Entscheidungssituation in vier Bestandteile.

\begin{itemize}
  \item Alternativen, Umweltzustände, Ergebnisse, Zielgröße (Nutzen, Gewinn)
  \item Matrixstruktur: Alternative × Umweltzustand → Ergebnis
\end{itemize}


\subsubsection{Informationsstand}


{\small
\begin{tabularx}{\linewidth}{XX}
\hline
\rowcolor{tableheadbg}
{\color{white}\textbf{Typ}} & {\color{white}\textbf{Bekannt}} \\[2pt]
\hline
\endhead
\hline
\endfoot
Sicherheit & Konsequenzen vollständig bekannt \\
Risiko & Wahrscheinlichkeiten bekannt \\
Unsicherheit & weder Wahrscheinlichkeiten noch Konsequenzen \\
\end{tabularx}
}



\subsection{18. Entscheidung bei Sicherheit}

\textbf{Beispiel Payment-System}: 200.000 Transaktionen/Jahr; Gesamtkosten = Fixkosten + Gebühr × Transaktionen.



{\small
\begin{tabularx}{\linewidth}{XXX}
\hline
\rowcolor{tableheadbg}
{\color{white}\textbf{Anbieter}} & {\color{white}\textbf{Gebühr/Transaktion}} & {\color{white}\textbf{Fixkosten/Jahr}} \\[2pt]
\hline
\endhead
\hline
\endfoot
A (Basic) & 0,50 € & 5.000 € \\
B (Pro) & 0,35 € & 20.000 € \\
C (Enterprise) & 0,25 € & 50.000 € \\
\end{tabularx}
}



\subsection{19. Entscheidung unter Risiko}

\textbf{Erwartungswert} = Σ (Ergebnis × Wahrscheinlichkeit) — Entscheidungskriterium bei bekannten Wahrscheinlichkeiten.



{\small
\begin{tabularx}{\linewidth}{XXX}
\hline
\rowcolor{tableheadbg}
{\color{white}\textbf{Alt.}} & {\color{white}\textbf{Erfolg (p=0,6)}} & {\color{white}\textbf{Misserfolg (p=0,4)}} \\[2pt]
\hline
\endhead
\hline
\endfoot
A & 100 & 0 \\
B & 200 & −50 \\
\end{tabularx}
}



\subsection{20. Entscheidung unter Unsicherheit}


{\small
\begin{tabularx}{\linewidth}{XXX}
\hline
\rowcolor{tableheadbg}
{\color{white}\textbf{Regel}} & {\color{white}\textbf{Logik}} & {\color{white}\textbf{Risikohaltung}} \\[2pt]
\hline
\endhead
\hline
\endfoot
Maximax (Optimismus) & bestes Ergebnis je Alternative → Maximum & risikofreudig \\
Maximin / Wald (Pessimismus) & schlechtestes Ergebnis je Alternative → Maximum & risikoavers \\
Minimum-Regret / Savage-Niehans & maximales Bedauern je Alternative → Minimum & Vermeidung Fehlentscheidung \\
Laplace (Gleichverteilung) & alle Zustände gleich wahrscheinlich → Durchschnitt & neutral \\
Hurwicz (Pessimismus-Optimismus) & Gewichtung Best-/Schlechtestfall (λ) & flexibel \\
\end{tabularx}
}



\subsection{21. Entscheidungen in der digitalen Welt}


{\small
\begin{tabularx}{\linewidth}{XX}
\hline
\rowcolor{tableheadbg}
{\color{white}\textbf{Klassisch}} & {\color{white}\textbf{Digital}} \\[2pt]
\hline
\endhead
\hline
\endfoot
wenig Daten & große Datenmengen \\
statische Modelle & kontinuierliche Optimierung \\
\end{tabularx}
}


\textbf{Beispiele}: A/B-Testing, Recommendation Systems, Dynamic Pricing.


\begin{quote}
Normative Modelle + Daten = bessere Entscheidungen.
\end{quote}




% ──────────────────────────────────────────────────────────────

\end{document}
